% Response to the reviewers for ct-2026-01838j (JCTC). Standalone document.
% Build from paper/ so that % Generated by tools/paper_numbers.py -- do not edit by hand.
% Every macro below is derived from a record in results/records/;
% tools/audit.py verifies that the derivation still reproduces.
\newcommand{\fsLOSkipLoPct}{82}
\newcommand{\fsLOSkipHiPct}{94}
\newcommand{\fsFreeAurocLo}{0.60}
\newcommand{\fsFreeAurocHi}{0.70}
\newcommand{\fsStableRankLo}{1.2}
\newcommand{\fsStableRankHi}{4.3}
\newcommand{\fsNaiveAccuracyRatio}{2.7}
\newcommand{\fsCovNCalMin}{84}
\newcommand{\fsCovNCalMax}{428}
\newcommand{\fsCovRealisedMin}{0.952}
\newcommand{\fsCovBetaMin}{0.953}
\newcommand{\fsCovSplits}{28}
\newcommand{\fsCovInfeasible}{0}
\newcommand{\fsFactorALo}{0.737}
\newcommand{\fsFactorAHi}{0.848}
\newcommand{\fsGlobalWins}{24}
\newcommand{\fsGlobalWinsSig}{19}
\newcommand{\fsGlobalCells}{24}
\newcommand{\fsFreeSig}{9}
\newcommand{\fsFreeSigTotal}{10}
\newcommand{\fsFreeGainLo}{0.067}
\newcommand{\fsFreeGainHi}{0.182}
\newcommand{\fsLOSkipLo}{0.822}
\newcommand{\fsLOSkipHi}{0.943}
\newcommand{\fsLOBeatsCV}{5}
\newcommand{\fsLOSystems}{6}
\newcommand{\fsBetaMAE}{0.015}
\newcommand{\fsMFourRetained}{0.970}
\newcommand{\fsMFourWorstSubset}{0.732}
\newcommand{\fsMTwoWorstSubset}{0.548}
\newcommand{\fsNaiveOverconfLo}{2.28}
\newcommand{\fsSharedOverconfHi}{25.77}
\newcommand{\fsOverconfNaiveSame}{2.58}
\newcommand{\fsOverconfSharedSame}{25.77}
\newcommand{\fsOverconfNaiveResampled}{2.56}
\newcommand{\fsOverconfSharedResampled}{13.52}
\newcommand{\fsOverconfNaiveWaterMD}{2.28}
\newcommand{\fsOverconfSharedWaterMD}{6.40}
\newcommand{\fsOverconfNaiveWaterPIMD}{3.08}
\newcommand{\fsOverconfSharedWaterPIMD}{24.08}
\newcommand{\fsNaivePassed}{3}
\newcommand{\fsMPtrajN}{136,923}
\newcommand{\fsMPtrajGlobal}{0.798}
\newcommand{\fsMPtrajFree}{0.943}
\newcommand{\fsMPtrajRandom}{0.499}
\newcommand{\fsMPtrajExtensivity}{2.50}
\newcommand{\fsMPtrajJaccard}{0.794}
\newcommand{\fsAlphaStar}{0.75}
\newcommand{\fsNormAurocZero}{0.797}
\newcommand{\fsNormAurocStar}{0.821}
\newcommand{\fsNormWins}{2}
\newcommand{\fsNormLosses}{1}
\newcommand{\fsCompFree}{0.9425}
\newcommand{\fsCompBoth}{0.9427}
\newcommand{\fsCompMarginal}{0.0002}
\newcommand{\fsCompFixedSig}{8}
\newcommand{\fsCompFixedN}{10}
\newcommand{\fsCompFixedLo}{0.022}
\newcommand{\fsCompFixedHi}{0.243}
\newcommand{\fsWaterSpeedupBOne}{1.53}
\newcommand{\fsWaterSpeedupBSixteen}{1.88}
\newcommand{\fsWaterSpeedupBSixteenHTwo}{1.85}
\newcommand{\fsWaterSpeedupBSixteenBlk}{1.89}
\newcommand{\fsWaterSpeedupSpread}{2.3}
\newcommand{\fsThreeBPASpeedupBSixteen}{1.10}
 resolves:
%   pdflatex response_to_reviewers && bibtex response_to_reviewers && pdflatex x2
% Red bracketed text is a \fsTODO placeholder for a number or location that the
% J10/J11/J12 experiments supply; tools/audit.py refuses a release while any
% remain.
\documentclass[11pt]{article}
\usepackage[margin=1in]{geometry}
\usepackage[T1]{fontenc}
\usepackage{xcolor}
\usepackage{booktabs}
\usepackage{natbib}
\usepackage{url}
\usepackage{hyperref}
\hypersetup{colorlinks=true,linkcolor=blue!60!black,citecolor=blue!60!black,urlcolor=blue!60!black}

\providecommand{\fsTODO}[1]{\textcolor{red}{[#1]}}
% Generated by tools/paper_numbers.py -- do not edit by hand.
% Every macro below is derived from a record in results/records/;
% tools/audit.py verifies that the derivation still reproduces.
\newcommand{\fsLOSkipLoPct}{82}
\newcommand{\fsLOSkipHiPct}{94}
\newcommand{\fsFreeAurocLo}{0.60}
\newcommand{\fsFreeAurocHi}{0.70}
\newcommand{\fsStableRankLo}{1.2}
\newcommand{\fsStableRankHi}{4.3}
\newcommand{\fsNaiveAccuracyRatio}{2.7}
\newcommand{\fsCovNCalMin}{84}
\newcommand{\fsCovNCalMax}{428}
\newcommand{\fsCovRealisedMin}{0.952}
\newcommand{\fsCovBetaMin}{0.953}
\newcommand{\fsCovSplits}{28}
\newcommand{\fsCovInfeasible}{0}
\newcommand{\fsFactorALo}{0.737}
\newcommand{\fsFactorAHi}{0.848}
\newcommand{\fsGlobalWins}{24}
\newcommand{\fsGlobalWinsSig}{19}
\newcommand{\fsGlobalCells}{24}
\newcommand{\fsFreeSig}{9}
\newcommand{\fsFreeSigTotal}{10}
\newcommand{\fsFreeGainLo}{0.067}
\newcommand{\fsFreeGainHi}{0.182}
\newcommand{\fsLOSkipLo}{0.822}
\newcommand{\fsLOSkipHi}{0.943}
\newcommand{\fsLOBeatsCV}{5}
\newcommand{\fsLOSystems}{6}
\newcommand{\fsBetaMAE}{0.015}
\newcommand{\fsMFourRetained}{0.970}
\newcommand{\fsMFourWorstSubset}{0.732}
\newcommand{\fsMTwoWorstSubset}{0.548}
\newcommand{\fsNaiveOverconfLo}{2.28}
\newcommand{\fsSharedOverconfHi}{25.77}
\newcommand{\fsOverconfNaiveSame}{2.58}
\newcommand{\fsOverconfSharedSame}{25.77}
\newcommand{\fsOverconfNaiveResampled}{2.56}
\newcommand{\fsOverconfSharedResampled}{13.52}
\newcommand{\fsOverconfNaiveWaterMD}{2.28}
\newcommand{\fsOverconfSharedWaterMD}{6.40}
\newcommand{\fsOverconfNaiveWaterPIMD}{3.08}
\newcommand{\fsOverconfSharedWaterPIMD}{24.08}
\newcommand{\fsNaivePassed}{3}
\newcommand{\fsMPtrajN}{136,923}
\newcommand{\fsMPtrajGlobal}{0.798}
\newcommand{\fsMPtrajFree}{0.943}
\newcommand{\fsMPtrajRandom}{0.499}
\newcommand{\fsMPtrajExtensivity}{2.50}
\newcommand{\fsMPtrajJaccard}{0.794}
\newcommand{\fsAlphaStar}{0.75}
\newcommand{\fsNormAurocZero}{0.797}
\newcommand{\fsNormAurocStar}{0.821}
\newcommand{\fsNormWins}{2}
\newcommand{\fsNormLosses}{1}
\newcommand{\fsCompFree}{0.9425}
\newcommand{\fsCompBoth}{0.9427}
\newcommand{\fsCompMarginal}{0.0002}
\newcommand{\fsCompFixedSig}{8}
\newcommand{\fsCompFixedN}{10}
\newcommand{\fsCompFixedLo}{0.022}
\newcommand{\fsCompFixedHi}{0.243}
\newcommand{\fsWaterSpeedupBOne}{1.53}
\newcommand{\fsWaterSpeedupBSixteen}{1.88}
\newcommand{\fsWaterSpeedupBSixteenHTwo}{1.85}
\newcommand{\fsWaterSpeedupBSixteenBlk}{1.89}
\newcommand{\fsWaterSpeedupSpread}{2.3}
\newcommand{\fsThreeBPASpeedupBSixteen}{1.10}


% A reviewer's point, quoted verbatim.
\newcommand{\rpoint}[1]{\begin{quote}\itshape #1\end{quote}}
\newcommand{\response}{\medskip\noindent\textbf{Response.} }
\newcommand{\changes}{\medskip\noindent\textbf{Changes.}}

\title{Response to the reviewers, manuscript ct-2026-01838j}
\author{Ian Poon}
\date{November 2026}

\begin{document}
\maketitle

\noindent Dear Prof.\ Liu,

\medskip\noindent
Thank you for the two reviews. Both reviewers found the submitted text hard to
read, found the committees and their accuracy insufficiently described, found
the literature coverage thin, and found the materials result claimed more
broadly than the evidence allowed. We agree on all four counts. The manuscript
has been rewritten from the first page: the bespoke vocabulary is gone, the
structure is conventional (Introduction, Methods, Results, Discussion), every
committee is described down to its loss function and training-data assignment,
a table of energy and force errors and a calibration figure have been added,
and the introduction now situates the work in the uncertainty-quantification
literature the reviewers pointed us to. Three new experiments address the
substance of the materials criticism: an out-of-distribution scan on OMat24
with all four foundation committees, a filtered re-analysis of the MPtraj pool
that removes near-equilibrium and single-atom structures, and one real
active-learning round on 3BPA with retraining. The claims have been rescoped to
what these experiments show. Below, each point is quoted in full and followed
by our answer and a list of the concrete changes. Two corrections we found
ourselves are listed at the end.

\section*{Reviewer 1}

\subsection*{Point 1.1}
\rpoint{The manuscript is difficult to follow. The large amount of
manuscript-specific terminology (``Factor A/B,'' ``oracle,'' ``preservation,''
``lane,'' ``gate,'' etc.) makes the main scientific questions unnecessarily
difficult to understand. I suggest substantial rewriting with clearer
definitions and a more straightforward presentation.}

\response
We agree. The manuscript has been rewritten in plain language with a
conventional structure. Every term the reviewer lists has been removed, from
the text, the captions and the tables, and replaced by a description of the
quantity itself. The two questions the paper asks are now stated as questions:
does committee disagreement identify the structures on which the model's
forces are most wrong (we call this error-ranking quality), and how faithfully
does a cheap screening rule reproduce the exact committee statistic. A
projection of the per-head force matrix is called what it costs, one backward
pass through the trunk. The rule that skips the exact computation is called a
screening rule. The two quantities that cost no backward pass are called
zero-cost baselines and named individually (the predicted force norm and the
spread of the head energies). Pronouns without a referent in the same sentence
and constructions of the form ``the former'' have been removed throughout.

\changes
\begin{itemize}
\item Full rewrite of all sections (Introduction, Methods, Results,
  Discussion); no section of the submitted text survives unchanged.
\item A notation table (Table~\ref{tab:notation}) at the end of Methods lists
  every symbol and every named score with its definition and cost.
\item Each structure-level score is defined at first use in Methods,
  Section~2.2 (``the sum over all $3N$ coordinates of the per-coordinate head
  variance'', ``max-component'', ``max-atom'').
\item Figure and table captions use the same names as the text.
\end{itemize}

\subsection*{Point 1.2}
\rpoint{The committee models are insufficiently characterized. Since the main
conclusions concern committee uncertainty, the manuscript should clearly
describe how the committee members/heads are trained, how their diversity is
generated, what loss function is used, and how the uncertainty is defined and
calibrated. This is particularly important given the large calibration
difference reported between the shared-trunk and independently trained
committees.}

\response
We agree, and this information was available but not written down. The new
Methods section gives the complete recipe of every committee. For the custom
shared-trunk committees: the MACE architecture (two interaction layers, 32
channels, maximum angular order one, correlation order three, cutoff
6.0~\AA), the loss (a weighted sum of energy and force mean-squared errors
with weights 10 and 1000, per-atom normalised as in MACE), the optimiser
(Adam with an exponential moving average and stochastic weight averaging over
the final quarter of training), and the head construction (each member is a
separate readout attached to every layer's features, all sharing the trunk).
Head diversity comes from the training-data assignment and is now stated
explicitly for each of the three constructions: disjoint heads (the training
set is split evenly across the eight heads with no overlap), overlapping
heads (each head sees an independently drawn 80\% of the training set), and
identical data (all heads see the full set and differ only by their random
initialisation). The independently trained committees are eight full models
of the same architecture with seeds zero to seven. The foundation committees
are eight new disjoint readout heads on a frozen MACE-MP-0b trunk, fine-tuned
on 8,000 (or 1,000) MPtraj structures, with the loss, optimiser and epoch
count stated. The uncertainty is defined exactly as it is computed: the
per-coordinate variance over heads, the per-atom standard deviation as the
mean over the three Cartesian components of its square root, and the three
structure-level scores built from these. Calibration is now assessed with two
standard metrics in addition to the overconfidence ratio: the variance of the
z-scores (error divided by predicted standard deviation, which is one for a
calibrated Gaussian) and reliability curves (binned error against binned
predicted standard deviation). The calibration gap the reviewer refers to is
now quantified with both: the shared-trunk committees have z-score variances
far above one (\fsTODO{range of z-score variance over shared-trunk
committees}), the independently trained committees are much closer to one
(\fsTODO{z-score variance of the independently trained committees}). The
likely cause is also discussed: heads that share a trunk share its errors, so
their spread underestimates the error of the shared part.

\changes
\begin{itemize}
\item Methods, Section~2.1 ``Committees'': the full recipe for each
  shared-trunk committee, each independently trained committee and each
  foundation committee, including the training-data assignment that
  generates head diversity, the loss function, the optimiser and the epoch
  counts, with a citation to \citet{beck2025committee} for the archive the
  checkpoints come from.
\item Methods, Section~2.2 ``Uncertainty definitions'': the per-coordinate
  variance, the per-atom standard deviation and the three structure-level
  scores, written as equations.
\item Methods, Section~2.5 ``Calibration metrics'': the overconfidence ratio,
  the z-score variance and the reliability curve, each defined.
\item Results, Section~3.1 and Table~\ref{tab:accuracy}: the calibration
  metrics for every committee.
\item Discussion: a paragraph on why shared-trunk committees are
  overconfident and what that implies for a screening rule that is calibrated
  to the committee rather than to the error.
\end{itemize}

\subsection*{Point 1.3}
\rpoint{The underlying model accuracy and uncertainty--error relationship
should be shown explicitly. Standard energy/force MAE or RMSE should be
reported for the relevant models and datasets. I would also strongly suggest
an error-versus-uncertainty scatter or calibration plot. AUROC alone does not
show the actual error distribution or how reliably the uncertainty criterion
correlates with model error.}

\response
We agree. A new table reports, for every committee and every test set, the
energy mean absolute error and root-mean-square error per atom, the force
mean absolute error and root-mean-square error of the committee mean, the
mean of the per-head force root-mean-square errors, the median per-atom
predicted standard deviation, the overconfidence ratio, the z-score variance
and the Spearman rank correlation between the per-atom predicted standard
deviation and the per-atom force error. A new figure shows, for
representative committees, the per-atom force error against the per-atom
predicted standard deviation as a binned reliability curve with the
calibrated diagonal, and the per-structure version alongside. The rank
correlations are moderate on the molecules (\fsTODO{range of per-atom
Spearman correlation on the molecular systems}) and lower on liquid water
(\fsTODO{range of per-atom Spearman correlation on water}), which is
consistent with the AUROC values and now lets the reader see the error
distribution behind them. AUROC is retained as the ranking metric because the
acquisition decision is a ranking, but it is no longer the only view.

\changes
\begin{itemize}
\item Results, Section~3.1 ``Accuracy and calibration'' and
  Table~\ref{tab:accuracy}: energy and force errors, calibration metrics and
  rank correlations for every system and construction, generated from
  \texttt{results/records/j12\_accuracy.json}.
\item Figure~\ref{fig:calibration}: per-atom error against predicted standard
  deviation (reliability curves) with the calibrated diagonal, for the
  shared-trunk and independently trained 3BPA committees and for water.
\item Methods, Section~2.4 ``Error targets and ranking metric'': the error
  targets are defined and the reason for using a rank metric is stated
  without claiming more for it.
\end{itemize}

\subsection*{Point 1.4}
\rpoint{The conclusion regarding the force-magnitude baseline appears too
broad. The strong performance of the free force-magnitude signal on MPtraj is
interesting, but it may reflect a model- or dataset-specific correlation
between large forces and large prediction errors. More evidence is needed
before drawing a general conclusion about the usefulness of committee
uncertainty. In addition, since no active-learning retraining is performed,
identifying high-error structures should be clearly distinguished from
demonstrating a better active-learning acquisition strategy.}

\response
We agree on both counts, and the revision changes the evidence as well as the
wording. On the first, the submitted text did not state that the MACE-MP-0b
trunk was itself pretrained on MPtraj, so the MPtraj pool is in-distribution
for the trunk and out-of-sample only for the readout heads; the revision
states this plainly. The correlation between large forces and large force
errors is a known effect, and we do not claim the predicted force norm as a
discovery: the claim is that it is a strong baseline that committee methods
must be scored against. To find out where that holds, three things were done.
First, a pre-registered out-of-distribution scan on three OMat24 validation
subsets (ab initio molecular dynamics from PBE at 1000~K in NVT and NPT, and
rattled structures), with all four foundation committees. OMat24 post-dates
MACE-MP-0 and is the out-of-distribution set that \citet{beck2025committee}
used. Second, a filtered re-analysis of the MPtraj pool that removes
structures whose largest reference force is below 0.05~eV/\AA\ and single-atom
cells, plus AUROC computed within deciles of the reference force magnitude,
so the near-equilibrium structures cannot carry the result. Third, the
conclusions are rescoped: the predicted force norm is reported as the baseline
in every setting, and the claim that the committee ``adds nothing'' is now
confined to the settings where the nested comparison shows that
(\fsTODO{which settings after J10 and the filtered re-analysis}). On the
second count, the submitted text identified high-error structures and did not
retrain, and the Discussion said so, but the title and the abstract did not.
The revision adds one real active-learning round on 3BPA: from a pool of 2,139
frames at 1200~K, each rule (the three committee scores, the two zero-cost
baselines and random) selects 50 frames, a single-head model is retrained on
the enlarged training set with three seeds, and the retrained models are
scored on the pool frames no rule selected and on the 600~K and 300~K test
sets. Everywhere else in the paper the wording now distinguishes
``identifies high-error structures'' from ``is a better acquisition rule''.

\changes
\begin{itemize}
\item Methods, Section~2.7 ``Datasets, splits and contamination statement'':
  the MACE-MP-0b trunk was pretrained on MPtraj; the pool is in-distribution
  for the trunk and out-of-sample only for the heads.
\item Methods, Section~2.8 ``Out-of-distribution and active-learning
  protocols'': the J10 and J11 protocols, registered before either ran.
\item Results, Section~3.6 and Table~\ref{tab:ood}: AUROC of every score and
  both zero-cost baselines on the three OMat24 subsets for all four foundation
  committees (\fsTODO{OMat24 AUROC values}); the filtered MPtraj re-analysis
  and the within-decile AUROC (\fsTODO{filtered and within-decile AUROC
  values}).
\item Results, Section~3.7 and Table~\ref{tab:al}: force errors of the
  retrained models after one acquisition round, by rule
  (\fsTODO{active-learning outcome}).
\item Abstract, Introduction and Discussion: the force-norm result is stated
  as a baseline, with the trunk-pretraining caveat, and limited to the
  settings tested; a limitations paragraph states that a single acquisition
  round on one molecule does not establish an acquisition strategy in general.
\end{itemize}

\subsection*{General remark}
\rpoint{Overall, the manuscript contains potentially interesting
observations, but I do not think the current analysis sufficiently supports
its broader conclusions. Substantial additional analysis and rewriting would
be needed.}

\response
Both have been done; the additional analysis is listed under Points~1.3 and
1.4, and the rewriting under Point~1.1.

\section*{Reviewer 2}

\subsection*{Summary of the manuscript}
\rpoint{The paper discusses the evaluation of committee uncertainties and
different strategies to further process them in workflows such as query by
committee. The key results are:
\begin{itemize}
\item A comparison against the actual error is important
\item The selection rule for structures in a query by committee workflow
  should be size independent
\item A structure-wide average of force uncertainty performs better than
  choosing the structure with the largest disagreement of a single atom or a
  single force component.
\item Given the additional cost associated with committee force
  disagreements, the relevant benchmark to compare these strategies against
  should not be limited to the random selection of structures, but also
  include essentially-free metrics.
\end{itemize}}

\response
We thank the reviewer for this summary, which is accurate, and we have used
it as a check on the rewritten abstract: a reader should be able to recover
these four statements from the new text without the old vocabulary. One
clarification on the third item: the structure-wide sum of the per-coordinate
head variances ranks error better than the maximum over components or atoms
only at fixed system size; when sizes vary the sum grows with the number of
atoms and must be normalised, which is the second item.

\subsection*{Point 2.1}
\rpoint{Some of the basic claims are - to my knowledge - not true, or at
least not sufficiently referenced. The author claims that a comparison
between the uncertainty metric and true error is uncommon in the field of
uncertainty quantification and instead cheap uncertainty metrics are measured
against more expensive ones. The author further claims that these uncertainty
metrics are strongly size dependent. Neither of these claims are supported by
the given references. In fact, there is only a single reference to another
uncertainty-focused work and neither of the two claims are supported by it.}

\response
The reviewer is right about the first claim, and it is withdrawn. The
uncertainty-quantification literature routinely compares an uncertainty
estimate with the true error; \citet{kahle2022quality}, \citet{tan2023single},
\citet{busk2021calibrated}, \citet{kellner2024shallow} and
\citet{beck2025committee} all do so, and the revision says so and cites them.
What the submitted text should have said, and the revision now says, is
narrower: the literature on cheap approximations of a committee statistic
(sketches, subsampled heads, low-rank estimators) validates the approximation
against the exact committee, and our contribution is to evaluate both links,
the committee against the error and the approximation against the committee,
in one study. The second claim is kept but is now supported. That
committee-based scores grow with system size unless normalised is documented
by \citet{smith2018less}, who divide by the square root of the number of
atoms; by \citet{kellner2024shallow}, who discuss how the energy uncertainty
scales with the number of atoms; and by \citet{lortaraprasert2026size}, who
study the question directly. Our own measurement of the effect on the MPtraj
pool is reported with these citations rather than as a novel observation.

\changes
\begin{itemize}
\item Introduction: the sentence claiming that comparison with the true error
  is uncommon is removed; the paragraph now cites the works above as
  examples of that comparison and states the narrower point about
  approximation studies.
\item Introduction and Methods, Section~2.2: the size dependence of the
  summed score is introduced with citations to \citet{smith2018less},
  \citet{kellner2024shallow} and \citet{lortaraprasert2026size}.
\item The bibliography has been expanded substantially (see Point~2.3).
\end{itemize}

\subsection*{Point 2.2}
\rpoint{Across most sections, the manuscript is incomprehensible. It uses
terminology that is completely uncommon in the field. In the methods section
it introduces methodology and concepts in very broad strokes and without any
references to more detailed work. It excessively uses pronouns and ``the
former'', ``others'' and similar constructions in a way that makes it
impossible to identify what is being talked about in each sentence.}

\response
We agree, and the answer is the same rewrite described under Point~1.1 of
Reviewer~1. In addition to the vocabulary change, the Methods section is now
written at the level of detail a reader needs to reproduce each step, and
each method is tied to the work it comes from: query by committee to
\citet{seung1992query} and \citet{krogh1995neural}; committees for
interatomic potentials to \citet{artrith2012high}, \citet{gastegger2017machine},
\citet{schran2020committee} and \citet{imbalzano2021uncertainty}; shallow and
multi-head ensembles to \citet{kellner2024shallow},
\citet{vinchurkar2025gnn}, \citet{schafer2026shallow} and
\citet{beck2025committee}; the trace estimator to \citet{meyer2021hutchpp};
the conformal constant to \citet{vovk2005algorithmic} and
\citet{lei2018distribution}; the calibration metrics to
\citet{busk2021calibrated}, \citet{pernot2022long} and \citet{tran2020methods};
and the block bootstrap to \citet{kunsch1989jackknife}. Every sentence of the
new text was checked for a pronoun or a ``the former'' whose referent is not
in the same sentence, and such sentences were rewritten.

\changes
\begin{itemize}
\item Methods, Sections~2.1 to 2.8, each with the references above at the
  point where the method is introduced.
\item Table~\ref{tab:notation}: notation table.
\item Removal of every ``the former'', ``the latter'', ``others'' and
  unanchored pronoun from the manuscript.
\end{itemize}

\subsection*{Point 2.3}
\rpoint{The manuscript lacks a proper introduction to the field and
references to highly relevant work. A manuscript working on the topic of
shallow ensembles/multi-head committees should at the very least cite papers
such as doi.org/10.1088/2632-2153/ad594a or doi.org/10.1088/2632-2153/ae0bf0,
as well as the review and papers contained in the review
doi.org/10.1515/revce-2024-0028 to evaluate the performance of uncertainty
estimates. The full list of relevant omitted publications is much longer and
up to the author to collect.}

\response
We agree. The three works the reviewer names are the shallow-ensemble paper of
\citet{kellner2024shallow}, the graph-network shallow-ensemble paper of
\citet{vinchurkar2025gnn}, and the review of \citet{dai2025review}; all three
are now cited and discussed, and the introduction has been rewritten to place
the work in that literature. The new introduction covers committee
uncertainty for interatomic potentials and its use in active learning
\citep{artrith2012high,gastegger2017machine,smith2018less,zhang2019active,schran2020committee,imbalzano2021uncertainty,zaverkin2021exploration,kulichenko2023uncertainty,chepkasov2026active},
alternatives to committees such as Gaussian-process and feature-space
estimates
\citep{podryabinkin2017active,janet2019quantitative,musil2019fast,vandermause2020flare,zhu2023fast},
shallow and multi-head ensembles
\citep{kellner2024shallow,vinchurkar2025gnn,schafer2026shallow,beck2025committee},
the evaluation and calibration of uncertainty estimates
\citep{tran2020methods,busk2021calibrated,kahle2022quality,pernot2022long,levi2022evaluating,tan2023single,heid2024spatially,dai2025review},
conformal and distribution-free calibration
\citep{vovk2005algorithmic,lei2018distribution,hu2022conformal}, the
propagation of uncertainty into observables
\citep{kellner2024shallow,bigi2024rigidity,schwalbekoda2025information}, and
uncertainty for foundation models on materials pools
\citep{deng2023chgnet,batatia2023foundation,barrosoluque2024omat24,riebesell2025matbench,bilbrey2025foundation}.
Deep ensembles in the wider machine-learning literature are cited through
\citet{lakshminarayanan2017simple}, and the open question of size scaling
through \citet{lortaraprasert2026size}. We do not claim the list is complete,
but it now reflects the state of the field.

\changes
\begin{itemize}
\item Introduction: rewritten as a field introduction of about a page and a
  half with the citations above; the previous ``Related work'' section is
  merged into it.
\item Bibliography: expanded with the works listed above.
\end{itemize}

\subsection*{Point 2.4 (Figures 1 and 2)}
\rpoint{Figures 1 and 2 compare 3BPA disjoint, 3BPA overlapping, 3BPA same,
ethanol, aspirin and azobenzene. Why are the first three system + committee
type and the last three just a system? What committee type was used? Why
change what is being compared mid-figure?}

\response
The reviewer is right that the labels were inconsistent. All six panels are
one system evaluated with one committee construction. The three 3BPA panels
are three constructions of the same system (disjoint heads, overlapping
heads, identical data). The three rMD17 panels are one jointly trained
committee with disjoint heads, trained on 50 frames of each of the ten rMD17
molecules, and evaluated on the held-out frames of ethanol, aspirin and
azobenzene. The figures now label every panel as ``system, construction'',
so the rMD17 panels read ``ethanol, disjoint heads'' and so on, and the
caption and Methods state that the three rMD17 panels share one committee.

\changes
\begin{itemize}
\item Figures~\ref{fig:factora} and \ref{fig:gate}: every panel labelled
  ``system, construction''.
\item Methods, Section~2.1: the rMD17 committee is described as one jointly
  trained committee evaluated on three molecules.
\end{itemize}

\subsection*{Point 2.5 (Table 1)}
\rpoint{Table 1 shows that the force magnitude is the best ``uncertainty
estimate'', beating every pooling of committee disagreement as well as the
energy disagreement. Given the experimental setup, this revolutionary result
does not hold up to scrutiny. It is stated that the model used was the
MACE-MP foundation model with disjoint committee heads. This model was trained
on the same MPtraj dataset used for testing. Furthermore, the dataset contains
many crystal structures at equilibrium, where forces are exactly 0. This has
massive implications on a test that uses force magnitude as an ``uncertainty
estimate''. Neither are these aspects discussed, nor are the results tested
on an out-of-domain test set.}

\response
We agree with the two factual criticisms and have acted on both. The
MACE-MP-0b trunk was pretrained on MPtraj; the submitted text did not say so,
and the revision states it plainly, together with the consequence that the
MPtraj pool is in-distribution for the trunk and out-of-sample only for the
readout heads. The pool does contain many near-equilibrium structures with
very small reference forces, and on such structures the predicted force norm
is small and the force error is small for a trivial reason. The revision
adds a filtered re-analysis that removes structures whose largest reference
force is below 0.05~eV/\AA\ and single-atom cells, and reports AUROC within
deciles of the reference force magnitude, so the ordering of the scores can
be read without that contribution. The revision also adds the
out-of-distribution test the reviewer asks for: a pre-registered scan on
three OMat24 validation subsets with all four foundation committees. OMat24
post-dates MACE-MP-0 and is the out-of-distribution set used by
\citet{beck2025committee}. We do not describe the force-norm result as
revolutionary, and the revision does not either. Force error grows with
force magnitude, which is well known, so the predicted force norm is a
strong baseline; the claim we make is that committee methods must be scored
against it, and the out-of-distribution, filtered and retraining evidence now
bounds where it is sufficient and where the committee adds information
(\fsTODO{summary of where the committee beats the force norm after J10 and
the filtered re-analysis}).

\changes
\begin{itemize}
\item Methods, Section~2.7: contamination statement for the MACE-MP-0b trunk.
\item Results, Section~3.6: filtered MPtraj re-analysis and within-decile
  AUROC (\fsTODO{values}); Table~\ref{tab:ood}: OMat24 results for all four
  foundation committees (\fsTODO{values}).
\item Table~\ref{tab:mptraj}: caption rewritten to state the trunk
  pretraining and to point to the filtered analysis; the word ``free'' is
  replaced by ``zero-cost baseline'' and the baseline is no longer described
  as beating the committee without qualification.
\item Abstract and Discussion: rescoped as described under Point~1.4 of
  Reviewer~1.
\end{itemize}

\subsection*{Concluding assessment}
\rpoint{In conclusion, the manuscript shows large problems in many core areas,
such as knowledge of the state of the art of the field, the execution of the
scientific experiments, the presentation of the data and the writing in
general. The key results of the manuscript beyond these issues are the
comparison of different pooling techniques for query by committee, the
comparison against essentially-free metrics and an analysis of the effects
of reducing the committee size of the uncertainty metric. These results of
some interest to the field of uncertainty quantification in atomistic
simulations. However, even after a hypothetical major revision that rectifies
all issues, the results would still not be of sufficient relevance to justify
a publication in a journal such as JCTC.}

\response
We have tried to rectify all of the issues, and we have also changed what the
paper offers. The three results the reviewer lists as being of some interest
are kept and are now stated clearly. The revision adds, on top of them, a
complete characterisation of the committees with their accuracy and
calibration, an out-of-distribution test of the zero-cost baseline on a
dataset the foundation model did not see, a filtered analysis that removes
the near-equilibrium confound, and a retraining experiment that tests whether
the ranking differences between rules survive one acquisition round. We
think that the resulting paper answers a question practitioners face
directly, namely when the backward passes of a force committee are worth
computing and what the alternative costs, and we leave the judgement of
relevance to the editor.

\section*{Corrections made on our own initiative}

\begin{enumerate}
\item \textbf{A mismatched comparison of overconfidence ratios.} The submitted
  text (Section ``Independently trained committees'') compared an
  overconfidence ratio of $\fsNaiveOverconfLo\times$ for an independently
  trained committee with $\fsSharedOverconfHi\times$ for ``the matched
  shared-trunk committee''. The first number belongs to the water
  molecular-dynamics committee and the second to the 3BPA identical-data
  committee; the two are not a matched pair. The revision reports the four
  matched pairs (3BPA identical data, 3BPA resampled data, water molecular
  dynamics, water path-integral molecular dynamics), each independently
  trained committee against the shared-trunk committee with the same training
  data and evaluation frames: 3BPA identical data
  $\fsOverconfNaiveSame\times$ against $\fsOverconfSharedSame\times$, 3BPA
  resampled data $\fsOverconfNaiveResampled\times$ against
  $\fsOverconfSharedResampled\times$, water molecular dynamics
  $\fsOverconfNaiveWaterMD\times$ against $\fsOverconfSharedWaterMD\times$,
  and water path-integral molecular dynamics
  $\fsOverconfNaiveWaterPIMD\times$ against
  $\fsOverconfSharedWaterPIMD\times$. The direction of the conclusion, that independently trained
  committees are far better calibrated in magnitude, is unchanged, and the
  revision now also reports the z-score variance for every pair.
\item \textbf{Wrong citation for MACE-MP.} The submitted text cited
  \citet{batatia2022mace} for the MACE-MP foundation model. That reference
  describes the MACE architecture; the foundation model is
  \citet{batatia2023foundation}. The revision cites each for what it is.
\end{enumerate}

\medskip\noindent
We are grateful to both reviewers for reviews that were detailed enough to act
on, and we hope the revised manuscript is now suitable for consideration.

\bibliographystyle{unsrtnat}
\bibliography{references}

\end{document}
